% s-153-p153-m002; original PDF 153, 154
\Aphorism{s-153-p153-m002}{241}
Dem feiner organisirten Manne begegnet es heute wohl oft genug, verstummen zu müssen. Von den Gleichbegabten, Ebenbürtigen geht Jeder seinen eigenen Weg; die rückwärtigen Beziehungen nach den niederen Stufen, seien es Individuen oder Zustände, sind abgebrochen — er findet sich allein, in einer Öde ohne Widerhall, jede seiner Kundgebungen, Wort, Werk, That, kehrt unbeantwortet zu ihm zurück. Kommt zu dieser Einsamkeit, diesem oft mehr als fünfjährigen Schweigen noch eine männliche Sensibilität hinzu, die sich in Augenblicksverbindungen auszugeben verschmäht, so ist mit Sicherheit vorauszusagen, dass es in nächster Zeit eine Katastrophe geben wird, nämlich bei der nächsten Begegnung mit einem Weibe. Das Weib ist es, in dem diese Aushungerung des Mittheilungstriebes, das zum Maximum angestaute Bedürfniss nach Umgang sich sättigt; das Bild des Weibes ist es, das die verschmachtenden Nerven mit Gier ergreifen, einsaugen und verwandeln. Kein Zweifel, diese Lösung des Knotens, diese Entspannung des gespannten Bogens ist ein Missverständniss, ein Nothbehelf, erfunden von der Ungeduld eines Kranken, eine verfrühte Krisis, der nicht Heilung, sondern neues Unheil folgt. Aller Aberglaube, zu jedem Individuum existire ein polar entsprechendes, eine ergänzende Hälfte im Sinne Platos, findet hier seine Nahrung; aber jedes andere Weib, das dem Vereinsamten zufällig als Erste begegnet wäre, hätte dieselbe Suggestion eingeflösst — man kommt, im Tumult einer solchen „Rettung“, nicht dazu, der Retterin ins Gesicht zu sehen. Hunger ist zwar der beste Koch, aber der schlechteste Kritiker und das gerade Gegentheil eines Feinschmeckers. — Nun, der Rausch verfliegt, und dem Enttäuschten tagt die Begrenztheit jedes Verständnisses, die Bedingtheit jedes persönlichen Verkehrs, die nothwendige Einschränkung des überquellenden Mittheilungsdranges: er sieht, dass zwischen Ich und Du ein kältehauchendes Vacuum klafft, das nur ein dünner Wärme- und Lichtstrahl kreuzt. Die durchgängige Mischung, das völlige Aufgehen eines Individuums im anderen ist Wahn und Romantik: Individuen sind incommensurabel und gehen nie in einander auf, es bleibt immer ein Rest, und nicht der schlechteste.
\par

